Typesetting is the craft of arranging type so that the reader forgets the type and sees only the text. A page that looks calm invites reading; a page that looks restless is tiring before the first word is read.

Two small accidents disturb the calm. A paragraph may start with a single line at the bottom of a page, and continue on the next one; this lonely line is called an orphan. Or a paragraph may end with a single line at the top of a page; this one is called a widow.

Both are easy to overlook while writing and hard to unsee once noticed. Typesetting systems therefore penalize them, and good ones repair them without asking the author to rewrite the text.
